% Glossary terms from ch20: lines of the form \item[Term] Definition.
\item[Trajectory prediction] Mapping the last $T_{\mathrm{obs}}$ observed positions of an agent (and possibly its neighbours) to its next $T_{\mathrm{pred}}$ positions, ideally as a distribution per horizon step.
\item[ADE / FDE / minADE$_K$] Average displacement error, the mean Euclidean error over the prediction horizon, and final displacement error, the error at the last step, with per-horizon versions $\mathrm{ADE}_h$, $\mathrm{FDE}_h$; $\mathrm{minADE}_K$ and $\mathrm{minFDE}_K$ score the best of $K$ sampled trajectories and are not comparable with the ADE/FDE of a deterministic predictor.
\item[Constant-velocity baseline] Extrapolation of the last observed position with a velocity estimated from the last $k$ intervals; exact on straight flight, with a noise error that grows like $h/k$, and the baseline every learned predictor must beat.
\item[Kalman extrapolation] Running the predict step of the constant-velocity Kalman filter $h$ times without updates; the optimal linear predictor under its noise model, with a covariance that grows with the horizon.
\item[LSTM cell] A recurrent unit whose cell state is updated by gated addition, $\vect{c}_t=\vect{f}_t\odot\vect{c}_{t-1}+\vect{i}_t\odot\tilde{\vect{c}}_t$, with forget, input and output gates in $(0,1)$ produced by logistic functions; the direct gradient path through $\vect{c}$ is $\diag(\vect{f}_t)$.
\item[Sequence-to-sequence predictor] An encoder LSTM that summarises the observed displacements and a decoder LSTM that rolls the future out one step at a time from the handed-over state, with a read-out per step.
\item[Agent-centred frame] The input normalisation in which the last observed position is the origin, the last heading is the $x$-axis and displacements are divided by a typical step length; it makes the predictor translation and rotation invariant.
\item[Teacher forcing] Feeding the decoder the true previous displacement during training instead of its own prediction; fast to train but subject to exposure bias, the mismatch between training inputs and free-running inputs.
\item[Gaussian negative log-likelihood] The loss $\sum_d[\log\sigma_d+(y_d-\mu_d)^2/(2\sigma_d^2)]$ for a read-out that predicts a mean and a standard deviation per step, minimised by the conditional mean and the conditional variance.
\item[Social pooling] The Social-LSTM construction that sums the hidden states of the neighbours in each cell of a grid centred on the agent and feeds the embedded grid to the agent's LSTM.
\item[Attention] The operation $\operatorname{softmax}(\mat{Q}\mat{K}\T/\sqrt{d_k})\mat{V}$ that returns for every query a convex combination of the values weighted by the similarity of the query to the keys; over time steps it replaces recurrence, over agents it models interaction, at cost $\Theta(n^2 d)$.
\item[Positional encoding] Sinusoids of geometrically spaced frequencies added to the token embeddings so that attention, which is otherwise a set operation, can tell the order of the time steps.
\item[Multimodal prediction] A prediction that puts probability mass on several distinct futures (mixture density outputs, sampled rollouts of a GAN or CVAE, ensembles), needed because a squared-error predictor converges to the conditional mean, which averages the modes.
\item[Time-indexed occupancy region] The sequence of $p$-probability ellipses $\mathcal{E}_h(p)$ of a predicted distribution, one per horizon step, from which the planners derive an inflated radius, a cost region or a chance-constraint margin.
\item[Calibration] Agreement between a nominal probability $p$ and the fraction of true positions that fall inside the predicted $p$-ellipse; learned predictors are often over-confident and need a per-horizon inflation factor fitted on validation data.
\item[Non-autoregressive decoder] A decoder that emits all $H$ predicted steps in one pass, from queries built out of the positional encodings of the future steps rather than out of its own previous outputs; it has no exposure bias, but no step can condition on the step before it.
